\documentclass[11pt,a4paper]{article}

\usepackage[utf8]{inputenc}
\usepackage[T1]{fontenc}
\usepackage[french]{babel}
\usepackage[a4paper,margin=2.1cm]{geometry}
\usepackage{amsmath,amssymb,amsthm,mathtools}
\usepackage{lmodern}
\usepackage{enumitem}
\usepackage{hyperref}
\usepackage{xcolor}
\usepackage{fancyhdr}
\usepackage{lastpage}
\usepackage{tikz}
\usepackage{stmaryrd}
\usepackage{array}
\usepackage{tabularx}
\usepackage{mathrsfs}

\hypersetup{
    colorlinks=true,
    linkcolor=black,
    urlcolor=blue
}

\setlength{\parindent}{0pt}
\setlength{\parskip}{0.5em}

\pagestyle{fancy}
\fancyhf{}
\fancyhead[L]{Corrigé détaillé --- DS Nombres complexes}
\fancyhead[R]{Prépa scientifique}
\fancyfoot[C]{\thepage/\pageref{LastPage}}

\newtheorem*{remarque}{Remarque}
\newtheorem*{methode}{Méthode}
\newtheorem*{idee}{Idée}
\newtheorem*{attention}{Attention}

\newcommand{\C}{\mathbb{C}}
\newcommand{\R}{\mathbb{R}}
\newcommand{\U}{\mathbb{U}}
\newcommand{\ii}{\mathrm{i}}
\newcommand{\e}{\mathrm{e}}
\newcommand{\bareme}[1]{\hfill{\small\textbf{[#1 pts]}}}
\newcommand{\ds}{\displaystyle}
\newcommand{\conj}[1]{\overline{#1}}
\newcommand{\im}{\operatorname{Im}}
\newcommand{\re}{\operatorname{Re}}

\begin{document}

\begin{center}
    {\LARGE \textbf{Corrigé détaillé du devoir surveillé}}\\[0.4em]
    {\Large \textbf{Nombres complexes}}\\[0.6em]
    \rule{0.9\linewidth}{0.6pt}
\end{center}

\section*{Exercice 1 --- Racines de l’unité, barycentre complexe et polygones réguliers}

On fixe un entier \(n\ge 2\) et
\[
\U_n=\{\omega\in\C\mid \omega^n=1\}.
\]

\subsection*{Partie A --- Premières propriétés des racines \(n\)-ièmes de l’unité}

\begin{enumerate}[label=\textbf{A.\arabic*.}]
    \item \textbf{Déterminer explicitement les éléments de \(\U_n\).}

    Soit \(\omega\in\U_n\). Alors
    \[
    \omega^n=1.
    \]
    En particulier \(\omega\neq 0\), donc \(\omega\) s’écrit sous forme exponentielle
    \[
    \omega=r\e^{\ii\theta},\qquad r>0.
    \]
    Alors
    \[
    \omega^n=r^n\e^{\ii n\theta}=1.
    \]
    Or \(1=\e^{2k\ii\pi}\) pour tout \(k\in\mathbb{Z}\). On en déduit
    \[
    r^n=1
    \qquad\text{et}\qquad
    n\theta=2k\pi.
    \]
    Comme \(r>0\), on a nécessairement \(r=1\), puis
    \[
    \theta=\frac{2k\pi}{n}.
    \]
    Ainsi tout élément de \(\U_n\) est de la forme
    \[
    \omega_k=\e^{2k\ii\pi/n},\qquad k\in\mathbb{Z}.
    \]

    Réciproquement, pour tout \(k\in\mathbb{Z}\),
    \[
    \left(\e^{2k\ii\pi/n}\right)^n=\e^{2k\ii\pi}=1,
    \]
    donc \(\e^{2k\ii\pi/n}\in\U_n\).

    Finalement,
    \[
    \boxed{\U_n=\left\{\e^{2k\ii\pi/n}\mid k\in\mathbb{Z}\right\}.}
    \]

    \item \textbf{Montrer que
    \[
    \U_n=\left\{\e^{2\ii k\pi/n}\mid k\in\llbracket 0,n-1\rrbracket\right\}.
    \]}

    On procède par \textbf{double inclusion}.

    Posons
    \[
    E=\left\{\e^{2\ii k\pi/n}\mid k\in\llbracket 0,n-1\rrbracket\right\}.
    \]

    \underline{Montrons que \(\U_n\subset E\).}

    Soit \(\omega\in\U_n\). D’après la question précédente, il existe \(m\in\mathbb{Z}\) tel que
    \[
    \omega=\e^{2m\ii\pi/n}.
    \]
    Par division euclidienne de \(m\) par \(n\), il existe \(q\in\mathbb{Z}\) et \(r\in\llbracket 0,n-1\rrbracket\) tels que
    \[
    m=qn+r.
    \]
    Alors
    \[
    \omega=\e^{2(qn+r)\ii\pi/n}
    =\e^{2q\ii\pi}\e^{2r\ii\pi/n}
    =\e^{2r\ii\pi/n}.
    \]
    Comme \(r\in\llbracket 0,n-1\rrbracket\), on a \(\omega\in E\).

    Ainsi
    \[
    \U_n\subset E.
    \]

    \underline{Montrons que \(E\subset \U_n\).}

    Soit \(\omega\in E\). Il existe donc \(k\in\llbracket 0,n-1\rrbracket\) tel que
    \[
    \omega=\e^{2k\ii\pi/n}.
    \]
    Alors
    \[
    \omega^n=\left(\e^{2k\ii\pi/n}\right)^n=\e^{2k\ii\pi}=1.
    \]
    Donc \(\omega\in\U_n\).

    Ainsi
    \[
    E\subset \U_n.
    \]

    Par double inclusion,
    \[
    \boxed{\U_n=\left\{\e^{2\ii k\pi/n}\mid k\in\llbracket 0,n-1\rrbracket\right\}.}
    \]

    \item \textbf{Montrer que \(\U_n\) est stable par produit et par passage à l’inverse.}

    Soient \(\omega,\omega'\in\U_n\). Alors
    \[
    \omega^n=1,\qquad (\omega')^n=1.
    \]
    Donc
    \[
    (\omega\omega')^n=\omega^n(\omega')^n=1.
    \]
    Ainsi \(\omega\omega'\in\U_n\).

    De plus, \(\omega^n=1\) implique \(\omega\neq 0\), donc \(\omega^{-1}\) existe, et
    \[
    (\omega^{-1})^n=(\omega^n)^{-1}=1.
    \]
    Donc \(\omega^{-1}\in\U_n\).

    Finalement, \(\U_n\) est stable par produit et par passage à l’inverse.

    \item \textbf{Calculer
    \[
    \sum_{\omega\in\U_n}\omega.
    \]}

    D’après A.2,
    \[
    \sum_{\omega\in\U_n}\omega=\sum_{k=0}^{n-1}\e^{2k\ii\pi/n}.
    \]
    Posons
    \[
    q=\e^{2\ii\pi/n}.
    \]
    Comme \(n\ge 2\), on a \(q\neq 1\), et la somme ci-dessus est géométrique :
    \[
    \sum_{k=0}^{n-1}q^k=\frac{1-q^n}{1-q}.
    \]
    Or
    \[
    q^n=\e^{2\ii\pi}=1.
    \]
    Donc
    \[
    \sum_{\omega\in\U_n}\omega=0.
    \]

    Ainsi
    \[
    \boxed{\sum_{\omega\in\U_n}\omega=0.}
    \]

    \item \textbf{Calculer
    \[
    S_m=\sum_{\omega\in\U_n}\omega^m
    \qquad (m\in\mathbb{Z}).
    \]}

    Écrivons les éléments de \(\U_n\) sous la forme
    \[
    \omega_k=\e^{2k\ii\pi/n},\qquad k=0,\dots,n-1.
    \]
    Alors
    \[
    S_m=\sum_{k=0}^{n-1}\left(\e^{2k\ii\pi/n}\right)^m
    =\sum_{k=0}^{n-1}\e^{2km\ii\pi/n}.
    \]
    C’est encore une somme géométrique, de raison
    \[
    q_m=\e^{2m\ii\pi/n}.
    \]

    \begin{itemize}
        \item Si \(q_m\neq 1\), alors
        \[
        S_m=\frac{1-q_m^n}{1-q_m}=0,
        \]
        car
        \[
        q_m^n=\e^{2m\ii\pi}=1.
        \]

        \item Si \(q_m=1\), alors chaque terme vaut \(1\), donc
        \[
        S_m=n.
        \]
    \end{itemize}

    Or
    \[
    q_m=1
    \Longleftrightarrow
    \e^{2m\ii\pi/n}=1
    \Longleftrightarrow
    n\mid m.
    \]

    Ainsi
    \[
    \boxed{
    S_m=
    \begin{cases}
    n & \text{si } n\mid m,\\
    0 & \text{sinon.}
    \end{cases}}
    \]

    \item \textbf{En déduire que, pour tout polynôme \(P\in\C_{n-1}[X]\),
    \[
    \frac1n\sum_{\omega\in\U_n}P(\omega)=P(0).
    \]}

    Soit
    \[
    P(X)=a_0+a_1X+\cdots+a_{n-1}X^{n-1}.
    \]
    Alors
    \[
    \sum_{\omega\in\U_n}P(\omega)
    =\sum_{\omega\in\U_n}\sum_{k=0}^{n-1}a_k\omega^k.
    \]
    Comme il s’agit de sommes finies, on peut intervertir les deux sommes :
    \[
    \sum_{\omega\in\U_n}P(\omega)
    =\sum_{k=0}^{n-1}a_k\sum_{\omega\in\U_n}\omega^k.
    \]
    D’après la question précédente,
    \[
    \sum_{\omega\in\U_n}\omega^0=n
    \qquad\text{et}\qquad
    \sum_{\omega\in\U_n}\omega^k=0\ \text{pour }1\le k\le n-1.
    \]
    Donc
    \[
    \sum_{\omega\in\U_n}P(\omega)=na_0.
    \]
    Or
    \[
    P(0)=a_0.
    \]
    Ainsi
    \[
    \boxed{\frac1n\sum_{\omega\in\U_n}P(\omega)=P(0).}
    \]
\end{enumerate}

\subsection*{Partie B --- Du cercle unité au polygone régulier}

On fixe \(a\in\C\), \(\rho>0\), et pour \(\omega\in\U_n\),
\[
z_\omega=a+\rho\omega.
\]

\begin{enumerate}[label=\textbf{B.\arabic*.}]
    \item \textbf{Montrer que les points d’affixes \(z_\omega\) appartiennent à un même cercle.}

    Soit \(\omega\in\U_n\). Comme \(\omega^n=1\), on a \(|\omega|=1\). Alors
    \[
    |z_\omega-a|=|\rho\omega|=\rho|\omega|=\rho.
    \]
    Donc tous les points d’affixes \(z_\omega\) sont à distance \(\rho\) du point d’affixe \(a\).

    Ils appartiennent donc tous au cercle de centre \(A(a)\) et de rayon \(\rho\).

    \item \textbf{Montrer que, si \(\omega\neq\omega'\),
    \[
    |z_\omega-z_{\omega'}|=\rho|\omega-\omega'|.
    \]}

    On calcule
    \[
    z_\omega-z_{\omega'}=(a+\rho\omega)-(a+\rho\omega')=\rho(\omega-\omega').
    \]
    En prenant les modules,
    \[
    |z_\omega-z_{\omega'}|=\rho|\omega-\omega'|.
    \]

    \item \textbf{En déduire que les points \(z_\omega\) sont les sommets d’un polygone régulier.}

    Les racines \(n\)-ièmes de l’unité
    \[
    \omega_k=\e^{2k\ii\pi/n},\qquad k=0,\dots,n-1,
    \]
    sont régulièrement réparties sur le cercle unité : elles ont toutes le même module \(1\), et deux racines consécutives correspondent à un incrément d’argument constant égal à \(2\pi/n\).

    L’application
    \[
    \omega\longmapsto a+\rho\omega
    \]
    est la composée d’une homothétie de rapport \(\rho\) et d’une translation. Elle conserve donc l’ordre cyclique et l’espacement angulaire.

    Les points \(z_\omega\) sont donc les sommets d’un polygone régulier à \(n\) côtés.

    \item \textbf{Montrer que
    \[
    \frac1n\sum_{\omega\in\U_n}z_\omega=a.
    \]}

    On écrit
    \[
    \sum_{\omega\in\U_n}z_\omega
    =\sum_{\omega\in\U_n}(a+\rho\omega)
    =na+\rho\sum_{\omega\in\U_n}\omega.
    \]
    Or, d’après A.4,
    \[
    \sum_{\omega\in\U_n}\omega=0.
    \]
    Donc
    \[
    \sum_{\omega\in\U_n}z_\omega=na.
    \]
    En divisant par \(n\),
    \[
    \boxed{\frac1n\sum_{\omega\in\U_n}z_\omega=a.}
    \]

    \item \textbf{Interprétation géométrique.}

    La quantité
    \[
    \frac1n\sum_{\omega\in\U_n}z_\omega
    \]
    est l’affixe du barycentre des sommets du polygone, pris avec des coefficients égaux.

    Le résultat signifie donc que le barycentre des sommets d’un polygone régulier est son centre.

    \item \textbf{Montrer que
    \[
    \prod_{\omega\in\U_n}(X-z_\omega)=(X-a)^n-\rho^n.
    \]}

    Posons
    \[
    Y=X-a.
    \]
    Alors
    \[
    X-z_\omega=X-a-\rho\omega=Y-\rho\omega.
    \]
    Donc
    \[
    \prod_{\omega\in\U_n}(X-z_\omega)=\prod_{\omega\in\U_n}(Y-\rho\omega).
    \]
    Pour chaque facteur,
    \[
    Y-\rho\omega=\rho\left(\frac{Y}{\rho}-\omega\right).
    \]
    D’où
    \[
    \prod_{\omega\in\U_n}(Y-\rho\omega)
    =\rho^n\prod_{\omega\in\U_n}\left(\frac{Y}{\rho}-\omega\right).
    \]
    Or les racines du polynôme \(T^n-1\) sont exactement les éléments de \(\U_n\), donc
    \[
    \prod_{\omega\in\U_n}(T-\omega)=T^n-1.
    \]
    En prenant \(T=\dfrac{Y}{\rho}\), on obtient
    \[
    \prod_{\omega\in\U_n}\left(\frac{Y}{\rho}-\omega\right)
    =\left(\frac{Y}{\rho}\right)^n-1.
    \]
    Ainsi
    \[
    \prod_{\omega\in\U_n}(Y-\rho\omega)
    =\rho^n\left(\frac{Y^n}{\rho^n}-1\right)=Y^n-\rho^n.
    \]
    En revenant à \(Y=X-a\), il vient
    \[
    \boxed{\prod_{\omega\in\U_n}(X-z_\omega)=(X-a)^n-\rho^n.}
    \]

    \item \textbf{Retrouver à partir de cette identité que les \(z_\omega\) sont les sommets d’un polygone régulier.}

    Les racines du polynôme
    \[
    (X-a)^n-\rho^n
    \]
    sont les solutions de
    \[
    (X-a)^n=\rho^n.
    \]
    Cela équivaut à
    \[
    X-a=\rho\omega
    \qquad\text{avec }\omega\in\U_n,
    \]
    c’est-à-dire
    \[
    X=a+\rho\omega.
    \]
    Les racines sont donc les images des racines \(n\)-ièmes de l’unité par l’application \(\omega\mapsto a+\rho\omega\).

    Elles sont donc régulièrement réparties sur le cercle de centre \(a\) et de rayon \(\rho\) : ce sont les sommets d’un polygone régulier.

\end{enumerate}

\subsection*{Partie C --- Réciproque partielle}

On suppose
\[
z_k-a=\lambda\e^{2k\ii\pi/n},
\qquad k=0,\dots,n-1,
\]
avec \(\lambda\in\C^\ast\).

\begin{enumerate}[label=\textbf{C.\arabic*.}]
    \item \textbf{Montrer que tous les points \(z_k\) sont sur un même cercle.}

    Pour tout \(k\),
    \[
    |z_k-a|=\left|\lambda\e^{2k\ii\pi/n}\right|=|\lambda|\cdot 1=|\lambda|.
    \]
    Donc tous les points d’affixes \(z_k\) sont à distance constante \(|\lambda|\) du point d’affixe \(a\).

    Ils appartiennent donc au cercle de centre \(a\) et de rayon \(|\lambda|\).

    \item \textbf{Montrer que
    \[
    \sum_{k=0}^{n-1}z_k=na.
    \]}

    En sommant l’égalité
    \[
    z_k=a+\lambda\e^{2k\ii\pi/n},
    \]
    on obtient
    \[
    \sum_{k=0}^{n-1}z_k
    =na+\lambda\sum_{k=0}^{n-1}\e^{2k\ii\pi/n}.
    \]
    La somme des racines \(n\)-ièmes de l’unité est nulle, donc
    \[
    \boxed{\sum_{k=0}^{n-1}z_k=na.}
    \]

    \item \textbf{Montrer que, pour tout entier \(m\),
    \[
    \sum_{k=0}^{n-1}(z_k-a)^m=
    \begin{cases}
    0 & \text{si } n\nmid m,\\
    n\lambda^m & \text{si } n\mid m.
    \end{cases}
    \]}

    On a
    \[
    z_k-a=\lambda\e^{2k\ii\pi/n}.
    \]
    Donc
    \[
    (z_k-a)^m=\lambda^m\e^{2km\ii\pi/n}.
    \]
    En sommant,
    \[
    \sum_{k=0}^{n-1}(z_k-a)^m
    =\lambda^m\sum_{k=0}^{n-1}\e^{2km\ii\pi/n}.
    \]
    La somme de droite a déjà été calculée en A.5. Ainsi
    \[
    \boxed{
    \sum_{k=0}^{n-1}(z_k-a)^m=
    \begin{cases}
    0 & \text{si } n\nmid m,\\
    n\lambda^m & \text{si } n\mid m.
    \end{cases}}
    \]

    \item \textbf{Montrer que, pour tout polynôme \(P\in\C_{n-1}[X]\),
    \[
    \frac1n\sum_{k=0}^{n-1}P(z_k)=P(a).
    \]}

    Posons
    \[
    Q(X)=P(a+\lambda X).
    \]
    Comme \(P\) est de degré au plus \(n-1\), \(Q\) aussi.

    Or
    \[
    z_k=a+\lambda\e^{2k\ii\pi/n},
    \]
    donc
    \[
    P(z_k)=Q\!\left(\e^{2k\ii\pi/n}\right).
    \]
    Par conséquent,
    \[
    \sum_{k=0}^{n-1}P(z_k)=\sum_{\omega\in\U_n}Q(\omega).
    \]
    En appliquant A.6 au polynôme \(Q\), on obtient
    \[
    \frac1n\sum_{\omega\in\U_n}Q(\omega)=Q(0).
    \]
    Donc
    \[
    \frac1n\sum_{k=0}^{n-1}P(z_k)=Q(0)=P(a+\lambda\cdot 0)=P(a).
    \]
    Finalement,
    \[
    \boxed{\frac1n\sum_{k=0}^{n-1}P(z_k)=P(a).}
    \]

    \item \textbf{Interprétation géométrique.}

    Cette propriété signifie que la moyenne des valeurs d’un polynôme de degré au plus \(n-1\) sur les sommets d’un polygone régulier est égale à sa valeur au centre.

    C’est une traduction algébrique de la forte symétrie du polygone régulier.

\end{enumerate}

\subsection*{Partie D --- Cas particuliers classiques}

\subsubsection*{1. Triangle équilatéral}

\begin{enumerate}[label=\textbf{D1.\arabic*.}]
    \item \textbf{Montrer que si \(z_1+z_2+z_3=0\) et si \(z_1,z_2,z_3\) ont même module, alors ils sont les sommets d’un triangle équilatéral.}

    Posons
    \[
    r=|z_1|=|z_2|=|z_3|.
    \]
    Comme les points sont distincts, on a \(r\neq 0\).

    Considérons
    \[
    u_k=\frac{z_k}{r}.
    \]
    Alors
    \[
    |u_k|=1
    \qquad\text{et}\qquad
    u_1+u_2+u_3=0.
    \]
    Les \(u_k\) appartiennent donc au cercle unité, et leur barycentre est l’origine.

    Divisons l’égalité \(u_1+u_2+u_3=0\) par \(u_1\neq 0\) :
    \[
    1+\frac{u_2}{u_1}+\frac{u_3}{u_1}=0.
    \]
    Posons
    \[
    v=\frac{u_2}{u_1},
    \qquad
    w=\frac{u_3}{u_1}.
    \]
    Alors \(|v|=|w|=1\) et
    \[
    1+v+w=0,
    \]
    donc
    \[
    w=-1-v.
    \]
    Comme \(|w|=1\), on a
    \[
    1=|w|=|-1-v|=|1+v|.
    \]
    En élevant au carré,
    \[
    1=|1+v|^2=(1+v)(1+\conj v)=2+v+\conj v=2+2\re(v).
    \]
    Donc
    \[
    \re(v)=-\frac12.
    \]
    Comme \(|v|=1\), cela impose
    \[
    v=\e^{2\ii\pi/3}
    \qquad\text{ou}\qquad
    v=\e^{-2\ii\pi/3}.
    \]
    On en déduit que les trois nombres \(u_1,u_2,u_3\) sont, à permutation près,
    \[
    \alpha,\quad \alpha\e^{2\ii\pi/3},\quad \alpha\e^{4\ii\pi/3}.
    \]
    Ils correspondent donc aux sommets d’un triangle équilatéral centré en \(0\).

    En multipliant par \(r\), on conserve la nature de la figure. Ainsi \(z_1,z_2,z_3\) sont les sommets d’un triangle équilatéral.

    \item \textbf{Montrer la réciproque suivante : si \(z_1,z_2,z_3\) sont les sommets d’un triangle équilatéral centré en \(0\), alors
    \[
    z_1+z_2+z_3=0.
    \]}

    Si le triangle est équilatéral et centré en \(0\), alors, quitte à renuméroter,
    \[
    z_1=r\e^{\ii\theta},\qquad
    z_2=r\e^{\ii(\theta+2\pi/3)},\qquad
    z_3=r\e^{\ii(\theta+4\pi/3)}.
    \]
    Donc
    \[
    z_1+z_2+z_3
    =r\e^{\ii\theta}\left(1+\e^{2\ii\pi/3}+\e^{4\ii\pi/3}\right).
    \]
    Or \(1,\e^{2\ii\pi/3},\e^{4\ii\pi/3}\) sont les racines cubiques de l’unité, donc leur somme vaut \(0\). Ainsi
    \[
    \boxed{z_1+z_2+z_3=0.}
    \]

    \item \textbf{Montrer qu’un triangle équilatéral quelconque admet, quitte à renuméroter ses sommets, une écriture de la forme
    \[
    z_1=a+r\e^{\ii\theta},\qquad
    z_2=a+r\e^{\ii(\theta+2\pi/3)},\qquad
    z_3=a+r\e^{\ii(\theta+4\pi/3)},
    \]
    avec \(a\in\C\), \(r>0\) et \(\theta\in\R\).}

    Soit un triangle équilatéral quelconque. Son centre est un certain point d’affixe \(a\). Les trois sommets sont tous à la même distance \(r>0\) de ce centre, et leurs directions sont séparées d’un angle de \(2\pi/3\).

    On peut donc écrire, quitte à renuméroter,
    \[
    z_1=a+r\e^{\ii\theta},\qquad
    z_2=a+r\e^{\ii(\theta+2\pi/3)},\qquad
    z_3=a+r\e^{\ii(\theta+4\pi/3)}.
    \]
    Réciproquement, une telle écriture décrit immédiatement trois points situés sur un même cercle, régulièrement espacés de \(2\pi/3\) : ils forment donc un triangle équilatéral.

\end{enumerate}

\subsubsection*{2. Carré}

\begin{enumerate}[label=\textbf{D2.\arabic*.}]
    \item \textbf{Montrer que si \(z_1+z_2+z_3+z_4=0\) et si les modules sont égaux, alors, après renumérotation éventuelle, ces points sont les sommets d’un rectangle centré en \(0\).}

Posons
\[
|z_1|=|z_2|=|z_3|=|z_4|=r.
\]
Comme les points sont distincts, on a \(r>0\).

Considérons les complexes
\[
u_k=\frac{z_k}{r}\qquad (k=1,2,3,4).
\]
Alors
\[
|u_1|=|u_2|=|u_3|=|u_4|=1
\qquad\text{et}\qquad
u_1+u_2+u_3+u_4=0.
\]
Les \(u_k\) sont donc quatre points distincts du cercle unité, dont la somme est nulle.

Quitte à renuméroter, on peut écrire
\[
u_1=\e^{\ii\alpha},\qquad u_2=\e^{\ii\beta},\qquad
u_3=\e^{\ii\gamma},\qquad u_4=\e^{\ii\delta}.
\]
En divisant l’égalité
\[
u_1+u_2+u_3+u_4=0
\]
par \(u_1\neq 0\), on obtient
\[
1+v+w+t=0,
\]
où
\[
v=\frac{u_2}{u_1},\qquad
w=\frac{u_3}{u_1},\qquad
t=\frac{u_4}{u_1}.
\]
On a encore
\[
|v|=|w|=|t|=1.
\]

On va montrer que, quitte à permuter \(v,w,t\), l’un d’eux est l’opposé d’un autre.  
De
\[
1+v+w+t=0
\]
on tire
\[
1+v=-(w+t).
\]
En prenant les modules au carré :
\[
|1+v|^2=|w+t|^2.
\]
Or, comme \(|v|=|w|=|t|=1\),
\[
|1+v|^2=(1+v)(1+\overline v)=2+v+\overline v=2+2\Re(v),
\]
et
\[
|w+t|^2=(w+t)(\overline w+\overline t)=2+w\overline t+t\overline w
=2+2\Re(w\overline t).
\]
Donc
\[
\Re(v)=\Re(w\overline t).
\]
Les deux complexes \(v\) et \(w\overline t\) sont de module \(1\). Deux complexes de module \(1\) ayant même partie réelle sont soit égaux, soit conjugués. Ainsi
\[
v=w\overline t
\qquad\text{ou}\qquad
v=\overline w\,t.
\]

Dans le premier cas,
\[
vt=w.
\]
En reportant dans \(1+v+w+t=0\), on obtient
\[
1+v+vt+t=(1+v)(1+t)=0.
\]
Donc
\[
v=-1 \qquad\text{ou}\qquad t=-1.
\]

Dans le second cas,
\[
vw=t,
\]
et de même
\[
1+v+w+vw=(1+v)(1+w)=0,
\]
d’où
\[
v=-1 \qquad\text{ou}\qquad w=-1.
\]

Dans tous les cas, l’un des trois nombres \(v,w,t\) vaut \(-1\). Cela signifie qu’un des \(u_2,u_3,u_4\) est égal à \(-u_1\). Quitte à renuméroter,
\[
u_3=-u_1.
\]
Comme
\[
u_1+u_2+u_3+u_4=0,
\]
on a alors
\[
u_2+u_4=0,
\]
donc
\[
u_4=-u_2.
\]

Ainsi, après renumérotation éventuelle,
\[
u_3=-u_1,\qquad u_4=-u_2.
\]
En revenant aux \(z_k\),
\[
z_3=-z_1,\qquad z_4=-z_2.
\]

Les diagonales \([z_1z_3]\) et \([z_2z_4]\) ont donc le même milieu, à savoir l’origine. Le quadrilatère formé est donc un parallélogramme centré en \(0\).

De plus, ses quatre sommets appartiennent tous à un même cercle, celui de centre \(0\) et de rayon \(r\). Or un parallélogramme inscrit dans un cercle est un rectangle.

On conclut donc qu’après renumérotation éventuelle, les points d’affixes \(z_1,z_2,z_3,z_4\) sont les sommets d’un rectangle centré en \(0\).

    \item \textbf{Donner une condition supplémentaire simple pour que ce rectangle soit un carré.}

    Un rectangle est un carré si, et seulement si, deux côtés consécutifs ont même longueur.

    Dans la configuration
    \[
    z_1,\ z_2,\ -z_1,\ -z_2,
    \]
    cela revient à demander que les directions de \(z_1\) et \(z_2\) diffèrent d’un angle droit. Autrement dit, les sommets consécutifs doivent être obtenus par une rotation d’angle \(\pm\pi/2\).

    \item \textbf{Traduire cette condition sous forme complexe.}

    Multiplier un complexe par \(\ii\) correspond à une rotation d’angle \(\pi/2\), et multiplier par \(-\ii\) à une rotation d’angle \(-\pi/2\). La condition cherchée s’écrit donc
    \[
    \boxed{z_2=\pm \ii z_1}
    \]
    après renumérotation convenable.

\end{enumerate}

\newpage

\section*{Exercice 2 --- Cercle unité et transformation homographique}

On suppose \(|u|=1\) et \(u\neq 1\), et on pose
\[
w=\frac{1+u}{1-u}.
\]

\begin{enumerate}[label=\textbf{\arabic*.}]
    \item \textbf{Montrer que \(\conj{w}=-w\).}

    On calcule
    \[
    \conj{w}=\frac{1+\conj{u}}{1-\conj{u}}.
    \]
    Comme \(|u|=1\), on a
    \[
    u\conj{u}=1
    \qquad\Longrightarrow\qquad
    \conj{u}=\frac1u.
    \]
    Donc
    \[
    \conj{w}
    =\frac{1+\frac1u}{1-\frac1u}
    =\frac{u+1}{u-1}
    =-\frac{1+u}{1-u}
    =-w.
    \]
    Ainsi
    \[
    \boxed{\conj{w}=-w.}
    \]

    \item \textbf{En déduire que \(w\in\ii\R\).}

    On sait qu’un complexe \(z\) est imaginaire pur si et seulement si
    \[
    \conj{z}=-z.
    \]
    Ici \(\conj{w}=-w\), donc
    \[
    \boxed{w\in \ii\R.}
    \]

    \item \textbf{Réciproquement, soit \(w\in\ii\R\). Montrer que
    \[
    u=\frac{w-1}{w+1}
    \]
    est bien défini et vérifie \(|u|=1\).}

    Écrivons
    \[
    w=\ii t,\qquad t\in\R.
    \]
    Alors \(w\neq -1\), puisque \(-1\notin\ii\R\), donc \(w+1\neq 0\), et \(u\) est bien défini.

    Ensuite,
    \[
    |u|=\left|\frac{w-1}{w+1}\right|=\frac{|w-1|}{|w+1|}.
    \]
    Or
    \[
    |w-1|^2=|-1+\ii t|^2=1+t^2,
    \qquad
    |w+1|^2=|1+\ii t|^2=1+t^2.
    \]
    Donc
    \[
    |w-1|=|w+1|.
    \]
    Il vient alors
    \[
    |u|=1.
    \]
    Ainsi
    \[
    \boxed{|u|=1.}
    \]

    \item \textbf{Montrer que
    \[
    u\longmapsto \frac{1+u}{1-u}
    \]
    réalise une bijection du cercle unité privé de \(1\) vers l’axe imaginaire pur.}

    On a déjà montré :
    \begin{itemize}
        \item si \(|u|=1\) et \(u\neq 1\), alors \(\dfrac{1+u}{1-u}\in\ii\R\) ;
        \item réciproquement, si \(w\in\ii\R\), alors
        \[
        u=\frac{w-1}{w+1}
        \]
        vérifie \(|u|=1\).
    \end{itemize}
    De plus, les deux formules sont réciproques : si
    \[
    w=\frac{1+u}{1-u},
    \]
    alors
    \[
    w(1-u)=1+u
    \Longrightarrow
    w-1=u(w+1)
    \Longrightarrow
    u=\frac{w-1}{w+1}.
    \]
    L’application est donc bijective.

    \item \textbf{Montrer que
    \[
    \frac{1+\e^{\ii\theta}}{1-\e^{\ii\theta}}
    =\ii \cot\left(\frac{\theta}{2}\right).
    \]}

    Posons \(u=\e^{\ii\theta}\). On factorise \(\e^{\ii\theta/2}\) :
    \[
    1+\e^{\ii\theta}
    =\e^{\ii\theta/2}\left(\e^{-\ii\theta/2}+\e^{\ii\theta/2}\right)
    =2\e^{\ii\theta/2}\cos\left(\frac{\theta}{2}\right).
    \]
    De même,
    \[
    1-\e^{\ii\theta}
    =\e^{\ii\theta/2}\left(\e^{-\ii\theta/2}-\e^{\ii\theta/2}\right)
    =-2\ii\,\e^{\ii\theta/2}\sin\left(\frac{\theta}{2}\right).
    \]
    Donc
    \[
    \frac{1+\e^{\ii\theta}}{1-\e^{\ii\theta}}
    =\frac{2\e^{\ii\theta/2}\cos(\theta/2)}{-2\ii\e^{\ii\theta/2}\sin(\theta/2)}
    =\ii\,\frac{\cos(\theta/2)}{\sin(\theta/2)}.
    \]
    Ainsi
    \[
    \boxed{
    \frac{1+\e^{\ii\theta}}{1-\e^{\ii\theta}}
    =\ii \cot\left(\frac{\theta}{2}\right).}
    \]

    \item \textbf{En déduire une expression de \(\cot(\theta/2)\) en fonction de \(\e^{\ii\theta}\).}

    En multipliant par \(-\ii\), on obtient
    \[
    \boxed{
    \cot\left(\frac{\theta}{2}\right)
    =-\ii\,\frac{1+\e^{\ii\theta}}{1-\e^{\ii\theta}}.}
    \]

    \item \textbf{Résoudre
    \[
    \frac{1+\e^{\ii\theta}}{1-\e^{\ii\theta}}=\ii\sqrt3
    \qquad\text{dans } [0,2\pi[.
    \]}

    Grâce à la question précédente,
    \[
    \ii\cot\left(\frac{\theta}{2}\right)=\ii\sqrt3.
    \]
    Donc
    \[
    \cot\left(\frac{\theta}{2}\right)=\sqrt3.
    \]
    Cela équivaut à
    \[
    \tan\left(\frac{\theta}{2}\right)=\frac1{\sqrt3}.
    \]
    On sait alors que
    \[
    \frac{\theta}{2}=\frac{\pi}{6}+k\pi,\qquad k\in\mathbb{Z}.
    \]
    Donc
    \[
    \theta=\frac{\pi}{3}+2k\pi.
    \]
    Dans \([0,2\pi[\), la seule solution est
    \[
    \boxed{\theta=\frac{\pi}{3}.}
    \]

    \item \textbf{Montrer que, pour \(z\neq 1\),
    \[
    |z|=1
    \Longleftrightarrow
    \frac{1+z}{1-z}\in\ii\R.
    \]}

    Le sens direct a déjà été démontré aux questions 1 et 2.

    Réciproquement, supposons
    \[
    \frac{1+z}{1-z}\in\ii\R.
    \]
    Posons
    \[
    w=\frac{1+z}{1-z}.
    \]
    Alors \(w\in\ii\R\), donc d’après la question 3,
    \[
    z=\frac{w-1}{w+1}
    \]
    vérifie \(|z|=1\).

    On a donc bien
    \[
    \boxed{|z|=1\Longleftrightarrow \frac{1+z}{1-z}\in\ii\R.}
    \]

    \item \textbf{Interpréter géométriquement ce résultat.}

    L’ensemble des points \(z\) tels que \(|z|=1\) est le cercle unité. La transformation
    \[
    z\longmapsto \frac{1+z}{1-z}
    \]
    envoie ce cercle privé du point \(1\) sur l’axe imaginaire pur.

    C’est donc un exemple de transformation homographique envoyant un cercle sur une droite.

\end{enumerate}

\newpage

\section*{Exercice 3 --- Un carré qui tourne}

Dans le plan complexe rapporté à un repère orthonormé direct, on considère deux points distincts \(A\) et \(B\) d’affixes respectives \(a\) et \(b\).

On cherche à construire les deux carrés ayant \([AB]\) pour côté.

\begin{enumerate}[label=\textbf{\arabic*.}]
    \item \textbf{Soit \(u=b-a\). Expliquer pourquoi multiplier \(u\) par \(\ii\) revient à effectuer une rotation d’angle \(\pi/2\).}

    Écrivons \(u\) sous forme exponentielle :
    \[
    u=r\e^{\ii\theta},\qquad r>0.
    \]
    Alors
    \[
    \ii u=\e^{\ii\pi/2}r\e^{\ii\theta}=r\e^{\ii(\theta+\pi/2)}.
    \]
    Le module est conservé, et l’argument est augmenté de \(\pi/2\). Multiplier par \(\ii\) correspond donc bien à une rotation d’angle \(\pi/2\).

    \item \textbf{Montrer que si \(C\) est le sommet d’un carré direct \(ABCD\), alors
    \[
    z_C-z_B=\ii(b-a).
    \]
    En déduire l’affixe de \(C\).}

    Dans un carré direct \(ABCD\), le vecteur \(\overrightarrow{BC}\) s’obtient en faisant subir au vecteur \(\overrightarrow{AB}\) une rotation d’angle \(\pi/2\), sans changer sa norme.

    Comme l’affixe de \(\overrightarrow{AB}\) est \(b-a\), celle de \(\overrightarrow{BC}\) est
    \[
    \ii(b-a).
    \]
    Donc
    \[
    z_C-z_B=\ii(b-a).
    \]
    D’où
    \[
    z_C=b+\ii(b-a).
    \]
    Ainsi
    \[
    \boxed{z_C=b+\ii(b-a).}
    \]

    \item \textbf{Montrer que l’autre carré construit sur \([AB]\) correspond à
    \[
    z_C-z_B=-\ii(b-a).
    \]
    Donner alors les deux affixes possibles de \(C\).}

    L’autre carré est obtenu en tournant \(\overrightarrow{AB}\) d’un angle \(-\pi/2\). L’affixe de \(\overrightarrow{BC}\) est alors
    \[
    -\ii(b-a).
    \]
    Donc
    \[
    z_C-z_B=-\ii(b-a),
    \]
    puis
    \[
    z_C=b-\ii(b-a).
    \]
    Les deux affixes possibles de \(C\) sont donc
    \[
    \boxed{z_C=b+\ii(b-a)\quad\text{ou}\quad z_C=b-\ii(b-a).}
    \]

    \item \textbf{En déduire les deux affixes possibles de \(D\).}

    Dans un carré \(ABCD\), le point \(D\) s’obtient à partir de \(A\) par la même rotation que \(C\) s’obtient à partir de \(B\). On a donc
    \[
    z_D-a=\pm\ii(b-a),
    \]
    avec le même signe que pour \(C\).

    Ainsi
    \[
    \boxed{z_D=a+\ii(b-a)\quad\text{ou}\quad z_D=a-\ii(b-a).}
    \]

    \item \textbf{Vérifier directement, dans chacun des deux cas, que les quatre côtés ont même longueur et que deux côtés consécutifs sont perpendiculaires.}

    Prenons par exemple le cas
    \[
    z_C=b+\ii(b-a),\qquad z_D=a+\ii(b-a).
    \]
    Alors
    \[
    z_B-z_A=b-a,
    \]
    \[
    z_C-z_B=\ii(b-a),
    \]
    \[
    z_D-z_C=(a+\ii(b-a))-(b+\ii(b-a))=a-b=-(b-a),
    \]
    \[
    z_A-z_D=a-(a+\ii(b-a))=-\ii(b-a).
    \]
    Les quatre côtés ont donc pour affixes
    \[
    b-a,\quad \ii(b-a),\quad -(b-a),\quad -\ii(b-a),
    \]
    qui ont tous le même module \(|b-a|\). Les quatre côtés ont donc même longueur.

    De plus, deux côtés consécutifs diffèrent par une multiplication par \(\ii\), donc sont perpendiculaires. Par exemple,
    \[
    z_C-z_B=\ii(z_B-z_A).
    \]
    Le quadrilatère est donc bien un carré.

    Le raisonnement est identique pour l’autre cas avec \(-\ii\).

    \item \textbf{On suppose maintenant
    \[
    a=1+\ii,\qquad b=4+2\ii.
    \]
    Calculer explicitement les affixes des deux points \(C\) possibles et des deux points \(D\) correspondants.}

    On a
    \[
    b-a=(4+2\ii)-(1+\ii)=3+\ii.
    \]
    Puis
    \[
    \ii(b-a)=\ii(3+\ii)=3\ii+\ii^2=-1+3\ii.
    \]
    Ainsi
    \[
    z_C^{(1)}=b+\ii(b-a)=(4+2\ii)+(-1+3\ii)=3+5\ii,
    \]
    \[
    z_D^{(1)}=a+\ii(b-a)=(1+\ii)+(-1+3\ii)=4\ii.
    \]

    De même,
    \[
    -\ii(b-a)=1-3\ii.
    \]
    Donc
    \[
    z_C^{(2)}=b-\ii(b-a)=(4+2\ii)+(1-3\ii)=5-\ii,
    \]
    \[
    z_D^{(2)}=a-\ii(b-a)=(1+\ii)+(1-3\ii)=2-2\ii.
    \]

    Finalement,
    \[
    \boxed{C_1: 3+5\ii,\quad D_1:4\ii,\qquad C_2:5-\ii,\quad D_2:2-2\ii.}
    \]

    \item \textbf{Déterminer l’affixe du centre \(G\) de chacun de ces carrés.}

    Le centre d’un carré est le milieu de l’une de ses diagonales, par exemple de \([AC]\). Donc
    \[
    z_G=\frac{a+z_C}{2}.
    \]

    Pour le premier carré :
    \[
    z_{G_1}=\frac{(1+\ii)+(3+5\ii)}{2}=\frac{4+6\ii}{2}=2+3\ii.
    \]

    Pour le second :
    \[
    z_{G_2}=\frac{(1+\ii)+(5-\ii)}{2}=\frac{6}{2}=3.
    \]

    Ainsi
    \[
    \boxed{z_{G_1}=2+3\ii,\qquad z_{G_2}=3.}
    \]

    \item \textbf{Montrer que les deux centres obtenus sont situés sur la médiatrice de \([AB]\).}

    Calculons d’abord le milieu \(M\) de \([AB]\) :
    \[
    z_M=\frac{a+b}{2}=\frac{(1+\ii)+(4+2\ii)}{2}=\frac{5+3\ii}{2}.
    \]
    La médiatrice de \([AB]\) est la droite passant par \(M\) et perpendiculaire à \(\overrightarrow{AB}\).

    Or les centres des deux carrés sont obtenus à partir de \(M\) en ajoutant ou retranchant un multiple de \(\ii(b-a)\) :
    \[
    z_{G_1}=z_M+\frac{\ii(b-a)}{2},\qquad
    z_{G_2}=z_M-\frac{\ii(b-a)}{2}.
    \]
    En effet,
    \[
    z_{G_1}=\frac{a+b+\ii(b-a)}{2},
    \qquad
    z_{G_2}=\frac{a+b-\ii(b-a)}{2}.
    \]
    Comme \(\ii(b-a)\) est perpendiculaire à \(b-a\), les vecteurs \(\overrightarrow{MG_1}\) et \(\overrightarrow{MG_2}\) sont perpendiculaires à \(\overrightarrow{AB}\).

    Donc \(G_1\) et \(G_2\) appartiennent à la médiatrice de \([AB]\).

    \item \textbf{Montrer qu’une rotation de centre \(a\) et d’angle \(\theta\) s’écrit
    \[
    z'-a=\e^{\ii\theta}(z-a).
    \]}

    Soit \(M\) un point d’affixe \(z\), et \(M'\) son image par la rotation de centre \(A(a)\) et d’angle \(\theta\).

    Le vecteur \(\overrightarrow{AM}\) a pour affixe \(z-a\). Après rotation d’angle \(\theta\), ce vecteur devient \(\overrightarrow{AM'}\), dont l’affixe s’obtient en multipliant par \(\e^{\ii\theta}\) :
    \[
    z'-a=\e^{\ii\theta}(z-a).
    \]
    C’est la formule générale d’une rotation de centre \(a\) et d’angle \(\theta\).

\end{enumerate}

\newpage

\section*{Exercice 4 --- Mais c’est qui ce nombre il est imaginaire ??}

On cherche \(z\in\C\setminus\{-1,1\}\) tel que :
\begin{itemize}
    \item \(|z|=1\),
    \item \(\dfrac{1+z}{1-z}\in\ii\R\),
    \item \(\dfrac{z-\ii}{z+\ii}\in\R\),
    \item \(\im(z)>0\).
\end{itemize}

\textbf{Déterminer \(z\).}

\bigskip

La première et la deuxième condition sont en fait liées. D’après l’exercice 2, on sait que
\[
|z|=1
\Longleftrightarrow
\frac{1+z}{1-z}\in\ii\R,
\qquad (z\neq 1).
\]
Ici, les deux premières informations reviennent donc à dire que \(z\) appartient au cercle unité, avec \(z\neq 1\).

La vraie nouvelle information est la suivante :
\[
\frac{z-\ii}{z+\ii}\in\R.
\]
Interprétons-la géométriquement.

\medskip

Si
\[
\frac{z-\ii}{z+\ii}\in\R,
\]
cela signifie que le quotient des deux complexes \(z-\ii\) et \(z+\ii\) est réel. Autrement dit, ces deux complexes ont le même argument modulo \(\pi\). Les vecteurs d’affixes \(z-\ii\) et \(z+\ii\) sont donc colinéaires.

Or :
\[
z-\ii = z-\ii,
\qquad
z+\ii = z-(-\ii).
\]
Ce sont respectivement les affixes des vecteurs
\[
\overrightarrow{I M}
\qquad\text{et}\qquad
\overrightarrow{J M},
\]
où \(I\) est le point d’affixe \(\ii\), \(J\) le point d’affixe \(-\ii\), et \(M\) le point d’affixe \(z\).

Dire que ces deux vecteurs sont colinéaires signifie que les points \(I\), \(J\), \(M\) sont alignés. Ainsi \(M\) appartient à la droite passant par \(\ii\) et \(-\ii\), c’est-à-dire à l’axe imaginaire.

Donc \(z\) est imaginaire pur :
\[
z=\ii y,\qquad y\in\R.
\]

Mais on sait aussi que \(|z|=1\), donc
\[
|\ii y|=|y|=1.
\]
Ainsi
\[
y=\pm 1,
\]
et donc
\[
z=\ii
\qquad\text{ou}\qquad
z=-\ii.
\]

Enfin, la condition
\[
\im(z)>0
\]
sélectionne le point situé dans le demi-plan supérieur. Parmi les deux possibilités, seule convient
\[
z=\ii.
\]

On conclut que
\[
\boxed{z=\ii.}
\]

\begin{remarque}
L’exercice était volontairement mis en scène comme une petite enquête : plusieurs indices semblaient pointer vers un complexe mystérieux, mais au final il s’agissait simplement de \( \ii \).
\end{remarque}

\vfill

\begin{center}
    \rule{0.9\linewidth}{0.6pt}\\[0.4em]
    \textit{Fin du corrigé}
\end{center}

\end{document}